% Add -5 times row 1 to row 3 of A, making the (3,1) entry zero; E=I3 with (3,1)=-5.
% Generated by generate.py; edit the generator and re-render.
\begin{tikzpicture}[x=1mm,y=1mm,font=\figurefont\fontsize{8pt}{9.6pt}\selectfont]
\path[use as bounding box] (-5,-5) rectangle (103,27);
\node[text=FigureMuted] at (-1.5,17) {$r_1$};
\node[text=FigureMuted] at (-1.5,11) {$r_2$};
\node[text=FigureMuted] at (-1.5,5) {$r_3$};
\node[text=FigureInk] at (9,24.5) {$\bm A$};
\draw[draw=FigureGrid,line width=0.4pt,fill=FigureBlueFill] (3,14) rectangle ++(6,6);
\node[text=FigureInk] at (6.0,17.0) {$1$};
\draw[draw=FigureGrid,line width=0.4pt,fill=FigureBlueFill] (9,14) rectangle ++(6,6);
\node[text=FigureInk] at (12.0,17.0) {$2$};
\draw[draw=FigureGrid,line width=0.4pt,fill=white] (3,8) rectangle ++(6,6);
\node[text=FigureInk] at (6.0,11.0) {$3$};
\draw[draw=FigureGrid,line width=0.4pt,fill=white] (9,8) rectangle ++(6,6);
\node[text=FigureInk] at (12.0,11.0) {$4$};
\draw[draw=FigureGrid,line width=0.4pt,fill=FigureYellowFill] (3,2) rectangle ++(6,6);
\node[text=FigureInk] at (6.0,5.0) {$5$};
\draw[draw=FigureGrid,line width=0.4pt,fill=FigureYellowFill] (9,2) rectangle ++(6,6);
\node[text=FigureInk] at (12.0,5.0) {$6$};
\draw[FigureStroke,line width=0.6pt] (4.2,20.6)--(2.4,20.6)--(2.4,1.4)--(4.2,1.4);
\draw[FigureStroke,line width=0.6pt] (13.8,20.6)--(15.6,20.6)--(15.6,1.4)--(13.8,1.4);
\draw[-{Stealth[length=2.2mm]},FigureStroke,line width=0.75pt] (19,11)--(52,11);
\node[text=FigureInk,above] at (35.5,11.5) {第3行加上第1行的$-5$倍};
\node[text=FigureInk,below] at (35.5,10.5) {$r_3\leftarrow r_3-5r_1$};
\node[text=FigureInk] at (61,24.5) {$\bm E\bm A$};
\draw[draw=FigureGrid,line width=0.4pt,fill=FigureBlueFill] (55,14) rectangle ++(6,6);
\node[text=FigureInk] at (58.0,17.0) {$1$};
\draw[draw=FigureGrid,line width=0.4pt,fill=FigureBlueFill] (61,14) rectangle ++(6,6);
\node[text=FigureInk] at (64.0,17.0) {$2$};
\draw[draw=FigureGrid,line width=0.4pt,fill=white] (55,8) rectangle ++(6,6);
\node[text=FigureInk] at (58.0,11.0) {$3$};
\draw[draw=FigureGrid,line width=0.4pt,fill=white] (61,8) rectangle ++(6,6);
\node[text=FigureInk] at (64.0,11.0) {$4$};
\draw[draw=FigureGrid,line width=0.4pt,fill=FigureGreenFill] (55,2) rectangle ++(6,6);
\node[text=FigureInk] at (58.0,5.0) {$0$};
\draw[draw=FigureGrid,line width=0.4pt,fill=FigureYellowFill] (61,2) rectangle ++(6,6);
\node[text=FigureInk] at (64.0,5.0) {$-4$};
\draw[FigureStroke,line width=0.6pt] (56.2,20.6)--(54.4,20.6)--(54.4,1.4)--(56.2,1.4);
\draw[FigureStroke,line width=0.6pt] (65.8,20.6)--(67.6,20.6)--(67.6,1.4)--(65.8,1.4);
\draw[FigureGrid,line width=0.6pt,dashed] (73.5,26)--(73.5,-2);
\node[text=FigureInk] at (89,24.5) {初等矩阵$\bm E$};
\draw[draw=FigureGrid,line width=0.4pt,fill=FigureBlueFill] (80,14) rectangle ++(6,6);
\node[text=FigureInk] at (83.0,17.0) {$1$};
\draw[draw=FigureGrid,line width=0.4pt,fill=FigureBlueFill] (86,14) rectangle ++(6,6);
\node[text=FigureInk] at (89.0,17.0) {$0$};
\draw[draw=FigureGrid,line width=0.4pt,fill=FigureBlueFill] (92,14) rectangle ++(6,6);
\node[text=FigureInk] at (95.0,17.0) {$0$};
\draw[draw=FigureGrid,line width=0.4pt,fill=white] (80,8) rectangle ++(6,6);
\node[text=FigureInk] at (83.0,11.0) {$0$};
\draw[draw=FigureGrid,line width=0.4pt,fill=white] (86,8) rectangle ++(6,6);
\node[text=FigureInk] at (89.0,11.0) {$1$};
\draw[draw=FigureGrid,line width=0.4pt,fill=white] (92,8) rectangle ++(6,6);
\node[text=FigureInk] at (95.0,11.0) {$0$};
\draw[draw=FigureGrid,line width=0.4pt,fill=FigureYellowFill] (80,2) rectangle ++(6,6);
\node[text=FigureInk] at (83.0,5.0) {$-5$};
\draw[draw=FigureGrid,line width=0.4pt,fill=white] (86,2) rectangle ++(6,6);
\node[text=FigureInk] at (89.0,5.0) {$0$};
\draw[draw=FigureGrid,line width=0.4pt,fill=white] (92,2) rectangle ++(6,6);
\node[text=FigureInk] at (95.0,5.0) {$1$};
\draw[FigureStroke,line width=0.6pt] (81.2,20.6)--(79.4,20.6)--(79.4,1.4)--(81.2,1.4);
\draw[FigureStroke,line width=0.6pt] (96.8,20.6)--(98.6,20.6)--(98.6,1.4)--(96.8,1.4);
\node[text=FigureMuted,font=\FigureNoteFont\fontsize{7.5pt}{9pt}\selectfont] at (89,-2.5) {对$\bm I_3$作同一变换};
\end{tikzpicture}
